\section{Appendix}
\label{sec:appendix}

\begin{table*}[!ht]
	\centering
	\begin{tabular}{lllllll}
		\toprule
		\textbf{Instance type}                    & \textbf{Accelerators} & \makecell{\textbf{Local SSD}\\\textbf{BW/GPU}} & \makecell{\textbf{Remote SSD}\\\textbf{BW/GPU}} & \makecell{\textbf{Network}\\\textbf{BW/GPU}} & \textbf{Has NVLink} & \textbf{Price} \\ \midrule
		a2-ultragpu-8g \cite{googleinstance}      & 8 x A100(80\,GB)      & 2.58\,Gbps                & 0.29\,Gbps                 & 12.5\,Gbps              & \ding{51}           & 40.44 USD/h    \\
		p4d.24xlarge\cite{awsinstance}            & 8 x A100(40\,GB)      & 2.31\,Gbps                & -                          & 100\,Gbps               & \ding{51}           & 45.039 USD/h   \\
		ml.hpcpni2.28xlarge\cite{volcanoinstance} & 8 x A100(80\,GB)      & 4\,Gbps                   & -                          & 100\,Gbps               & \ding{55}           & 48.23 USD/h    \\
		p4de.24xlarge\cite{awsinstance}           & 8 x A100(80\,GB)      & 2.31\,Gbps                & -                          & 100\,Gbps               & \ding{51}           & 56.328 USD/h   \\
		a3-highgpu-8g\cite{googleinstance}        & 8 x H100              & 6.09\,Gbps                & 0.97\,Gbps                 & 100\,Gbps               & \ding{51}           & 88.25 USD/h    \\
		a3-megagpu-8g\cite{googleinstance}        & 8 x H100              & 6.09\,Gbps                & 0.97\,Gbps                 & 200\,Gbps               & \ding{51}           & Unavailable    \\
		p5.48xlarge \cite{awsinstance}            & 8 x H100              & 9.8\,Gbps                 & -                          & 400\,Gbps               & \ding{51}           & Unavailable    \\ \bottomrule
	\end{tabular} \\[10pt]

    \begin{minipage}{1\linewidth}
        \caption{\small{\textit{
		A survey of {\maas} hardware configurations from GPU vendors.
        }}}
    \label{tab:vendor-hardware}
    \end{minipage} %\\[-10pt]	
\end{table*}

\subsection{Notable implementation details}
\label{sec:appendix-implementation}


\nospacestitle{Network library }
We implement a communication library
that abstracts both NVLink and RDMA to holistically transfer parameters,
similar to DeepEP~\cite{deepep}.
During our implementation, we found 
establishing communication group between machines is slow (e.g., 100\,ms) 
when using off-the-shelf group communicators (e.g., NCCL~\cite{nccl}),
which significantly limit the effectiveness of network-based scaling.
Fortunately, we found that our plan only requires P2P communication between each pair of nodes. 
Therefore, we pre-create a connection pool that supports full-mesh connections on each.
While the compute network (RDMA) has potential scalability issue~\cite{fasst}, 
it only occurs when transferring small payloads and can be addressed using advanced RDMA 
transport like DCT~\cite{krcore}. 

\stitle{Native serving engine with CUDA context pool. }
Before execution, a CUDA context with loaded kernels (cuModule) must be created on GPU.
Creating such a CUDA context takes about 500\,ms, and is non-negligible in serving instance
autoscaling. To mitigate such an overhead, {\sys} preserves a small CUDA context pool with
pre-loaded kernels and transfers parameters to GPU within one of the existing CUDA contexts,
similar to an existing work~\cite{DBLP:journals/corr/abs-2405-12079}.
Furthermore, {\sys} is built using C++ and native CUDA APIs, eliminating the overhead of initializing
PyTorch (e.g. \texttt{dlopen}).

\stitle{Fault tolerance. }
When machine failures occur, 
we will autoscale new instances using our scaling mechanism. 
One problem is that cached parameters on the failed machine are lost, 
so we need to redistribute these parameters to other machines to maintain 
our global parameter pool invariant. 
For other components in the system like 
scheduler or monitor failures, we follow the same procedure
in existing work for recovery~\cite{DBLP:conf/usenix/Romero0YK21,DBLP:journals/corr/abs-2401-14351}. 



\subsection{Hardware configurations for {\maas}}
\label{sec:appendix-hardware}

\noindent
Table~\ref{tab:vendor-hardware} lists the hardware configurations backed by typical 
{\maas} systems. 